\documentclass[11pt,a4paper]{article}
\usepackage[margin=1.8cm]{geometry}
\usepackage{amsmath,amssymb,tabularx,array,needspace,fontspec,cancel}
\setmainfont{FreeSerif}
\newfontfamily\urdufont[Script=Arabic]{FreeSerif}
\TeXXeTstate=1
\newcommand{\urdu}[1]{{\urdufont\beginR #1\endR}}
\newcommand{\lr}[1]{\beginL #1\endL}
\setlength{\parindent}{0pt}
\setlength{\parskip}{0pt}
\renewcommand{\arraystretch}{1.15}
% \umcq{n}{urdu question}{A}{B}{C}{D}; options shown D C B A left-to-right like the original
\newcommand{\umcq}[6]{%
  \Needspace{4\baselineskip}%
  \par\medskip{\urdufont\raggedleft\beginR\textbf{\lr{#1.}}\ #2\endR\par}\nopagebreak
  \noindent\begin{tabularx}{\linewidth}{@{}*{4}{>{\raggedleft\arraybackslash}X}@{}}
    \lr{(D)}~#6 & \lr{(C)}~#5 & \lr{(B)}~#4 & \lr{(A)}~#3
  \end{tabularx}\par}

\begin{document}
\begin{center}
  {\Huge\bfseries The Hope Science Academy}\\[3pt]
  {\small Excellence is our standard not goal}\\[2pt]
  {\small 247-E-1, (Opposite LDA School) Johar Town, Lahore. PH: 0300-8194789, 0322-4157001}
\end{center}
\medskip\noindent
\begin{tabularx}{\linewidth}{|X|X|X|X|}
\hline
\textbf{Student Name:} & \textbf{Roll Num:} & \textbf{Class:} 10th & \textbf{Paper Code:} \\ \hline
\textbf{Subject Name:} Mathematics. & \textbf{Time Allowed:} & \textbf{Total Marks:} 220 & \textbf{Exam Date:} 09-Oct-2026 \\ \hline
\multicolumn{4}{|l|}{\textbf{Exam Syllabus:} FULL BOOK} \\ \hline
\end{tabularx}
\bigskip
\noindent\textbf{1$\times$110=110}\hfill{\urdufont\beginR\textbf{سوال نمبر \lr{2}۔ چار ممکنہ جوابات میں سے درست پر دائرہ لگائیں}\endR}
\medskip
\umcq{1}{\lr{$i^{2}+i^{4}=$}}{-1}{0}{1}{2}
\umcq{2}{\lr{$(2-3i)(2+3i)$} کا حقیقی حصہ ہے:}{-3}{1}{4}{13}
\umcq{3}{\lr{$(2-i)(2+i)$} کا خیالاتی حصہ ہے:}{0}{1}{7}{9}
\umcq{4}{\lr{$x+iy$} کا ایک خالص خیالاتی عدد ہو گا، جب:}{y=0}{x=0}{\lr{$i=0$}}{\lr{x=0, y=0}}
\umcq{5}{\lr{$5-2i$} کا جمعی معکوس ہے:}{\lr{$5+2i$}}{\lr{$-5-2i$}}{\lr{$5-2i$}}{\lr{$-5+2i$}}
\umcq{6}{\lr{$z=1+i$} کا ضربی معکوس ہے؟:}{\lr{$1-i$}}{\lr{$i$}}{\lr{$\dfrac{1}{2}-\dfrac{1}{2}i$}}{\lr{$-\dfrac{1}{\sqrt{2}}+\dfrac{1}{\sqrt{2}}i$}}
\umcq{7}{اگر \lr{$z=4-3i$} ہو تو \lr{$z\bar{z}$} برابر ہے:}{3}{9}{16}{25}
\umcq{8}{\lr{$9-4i$} کا کانجوگیٹ ہے:}{\lr{$-9-4i$}}{\lr{$9+4i$}}{\lr{$9+9i$}}{\lr{$4-9i$}}
\umcq{9}{اگر \lr{$z=4+4i$} ہو تو \lr{$z+\bar{z}$} برابر ہے:}{8}{\lr{$8+8i$}}{\lr{$8i$}}{0}
\umcq{10}{اگر \lr{$z=5+4i$} ہو تو \lr{$|z|$} برابر ہے:}{9}{25}{41}{\lr{$\sqrt{41}$}}
\umcq{11}{مساوات \lr{$2x^{2}-x+1=0$} کی قسم ہے:}{\urdu{دو درجی}}{\urdu{یک درجی}}{\urdu{سہ درجی}}{\urdu{خالص دو درجی}}
\umcq{12}{\lr{$x^{2}+5x-5=0$} کا فرق کنندہ کیا ہے؟}{-20}{20}{25}{45}
\umcq{13}{\lr{$3x^{2}-9=0$} کا حل سیٹ ہے:}{\lr{\{3\}}}{\lr{\{$\pm3$\}}}{\lr{\{$\pm\sqrt{3}$\}}}{\lr{\{$\sqrt{3}$\}}}
\umcq{14}{\lr{$3x^{2}+5x-12=0$} کے روٹس کا مجموعہ ہے:}{\lr{$\dfrac{5}{3}$}}{\lr{$-\dfrac{5}{3}$}}{\lr{$\dfrac{12}{3}$}}{\lr{$\dfrac{3}{5}$}}
\umcq{15}{\lr{$3x^{2}+5x-12=0$} کے روٹس کا حاصل ضرب ہے:}{-4}{3}{4}{5}
\umcq{16}{\lr{$(x-3)(x+3)=0$} کے روٹس ہیں:}{\lr{3,-3}}{\lr{3,3}}{\lr{-3,-3}}{\lr{9,0}}
\umcq{17}{\lr{2} اور \lr{3} ...... کے روٹس ہیں:}{\lr{$x^{2}+5x+6=0$}}{\lr{$x^{2}+6x+5=0$}}{\lr{$x^{2}-5x+6=0$}}{\lr{$x^{2}+6x-5=0$}}
\umcq{18}{اگر \lr{$b^{2}-4ac>0$} اور مکمل مربع ہو تو \lr{$ax^{2}+bx+c=0$} کے روٹس ہیں:}{\urdu{مساوی}}{\urdu{غیر مساوی}}{\urdu{خیالاتی}}{\urdu{غیر ناطق}}
\umcq{19}{اگر \lr{$b^{2}-4ac=0$} ہو تو \lr{$ax^{2}+bx+c=0$} کے روٹس ہیں؟}{\urdu{غیر مساوی}}{\urdu{غیر ناطق}}{\urdu{خیالاتی}}{\urdu{مساوی}}
\umcq{20}{\lr{$x-2c=b$} کا سبجیکٹ \lr{c} ہے:}{\lr{x+b}}{\lr{b-x}}{\lr{$\dfrac{x-b}{2}$}}{\lr{$\dfrac{b-x}{2}$}}
\umcq{21}{اگر \lr{$\begin{bmatrix}a+2\\3\end{bmatrix}=\begin{bmatrix}7\\3\end{bmatrix}$} ہو تو \lr{a} برابر ہے:}{3}{5}{6}{7}
\umcq{22}{\lr{$A=\begin{bmatrix}3&5&0\end{bmatrix}$} ایک ...... قالب ہے:}{\urdu{قطاری}}{\urdu{مربعی}}{\urdu{کالمی}}{\urdu{صفری}}
\umcq{23}{\lr{$B=\begin{bmatrix}2\\3\end{bmatrix}$} ایک ...... قالب ہے:}{\urdu{وحدانی}}{\urdu{مربعی}}{\urdu{قطاری}}{\urdu{کالمی}}
\umcq{24}{\lr{$M=\begin{bmatrix}1&0&0\\0&1&0\\0&0&1\end{bmatrix}$} ایک ...... قالب ہے:}{\urdu{مستطیلی}}{\urdu{وحدانی}}{\urdu{کالمی}}{\urdu{قطاری}}
\umcq{25}{اگر \lr{$A^{t}=-A$} ہو تو \lr{A} ایک ...... قالب ہے:}{\urdu{سمیٹرک}}{\urdu{قطاری}}{\urdu{مستطیلی}}{\urdu{سکیو سمیٹرک}}
\umcq{26}{اگر \lr{$A=\begin{bmatrix}3&4\end{bmatrix}$, $B=\begin{bmatrix}7&8\end{bmatrix}$} ہو تو \lr{$A+B$} برابر ہے:}{\lr{[32 21]}}{\lr{[28 24]}}{\lr{[12 10]}}{\lr{[11 11]}}
\umcq{27}{اگر \lr{$A=\begin{bmatrix}3&4\end{bmatrix}$, $B=\begin{bmatrix}7&8\end{bmatrix}$} ہوں تو \lr{$B-A$} برابر ہے:}{\lr{[12 10]}}{\lr{[11 11]}}{\lr{[4 4]}}{\lr{[-4 -4]}}
\umcq{28}{\lr{$\begin{bmatrix}3\\4\end{bmatrix}$} کا جمعی معکوس کیا ہے؟}{\lr{$\begin{bmatrix}-3\\4\end{bmatrix}$}}{\lr{$\begin{bmatrix}-3\\-4\end{bmatrix}$}}{\lr{$\begin{bmatrix}3\\4\end{bmatrix}$}}{\lr{$\begin{bmatrix}-1\\1\end{bmatrix}$}}
\umcq{29}{اگر \lr{$A=\begin{bmatrix}3&1&3\\2&0&2\end{bmatrix}$} ہو تو \lr{$A^{t}$} کا مرتبہ ہے:}{\lr{3-by-2}}{\lr{2-by-3}}{\lr{3-by-3}}{\lr{2-by-2}}
\umcq{30}{\lr{$\begin{vmatrix}3&1\\0&4\end{vmatrix}=$} برابر ہے:}{11}{12}{13}{-11}
\umcq{31}{اگر \lr{$f(x)=\dfrac{5x-6}{3}$} ہو تو \lr{$f(3)$} کی قیمت کیا ہو گی؟}{-1}{3}{9}{15}
\umcq{32}{\lr{X} سے \lr{Y} تفاعل \lr{f} کو ...... سے ظاہر کیا جاتا ہے:}{\lr{$f:XY$}}{\lr{$f:Y\rightarrow X$}}{\lr{$f:X\rightarrow Y$}}{\lr{$f:\dfrac{X}{Y}$}}
\umcq{33}{\lr{$(fog)(x)=$}}{\lr{$(f+g)(x)$}}{\lr{$(f-g)(x)$}}{\lr{$f(g(x))$}}{\lr{$f(x)\div g(x)$}}
\umcq{34}{اگر \lr{$f(x)=2x+3,\ g(x)=x+1$} ہوں تو \lr{$f(x)+g(x)=$}:}{\lr{3x}}{\lr{3x+4}}{4}{\lr{$2x^{2}+3$}}
\umcq{35}{اگر \lr{$f(x)=5x+2,\ h(x)=2x-2$} ہوں تو \lr{$f(x)-h(x)=$}:}{\lr{3x}}{\lr{$5x^{2}-4$}}{\lr{3x+4}}{\lr{-3x-4}}
\umcq{36}{اگر \lr{$f(x)=3x+1,\ g(x)=2x$} ہوں تو \lr{$g(x)\times f(x)=$}:}{\lr{6x+2x}}{\lr{$5x^{2}+1$}}{\lr{x+1}}{\lr{$6x^{2}+2x$}}
\umcq{37}{اگر \lr{$f(x)=x^{2}-4,\ g(x)=x+2,\ x\neq-2$} ہوں تو \lr{$\dfrac{f(x)}{g(x)}=$}:}{\lr{$\dfrac{1}{x-2}$}}{\lr{$\dfrac{1}{x+2}$}}{\lr{x+2}}{\lr{x-2}}
\umcq{38}{مطلق قدر والے تفاعل کے گراف کی شکل کیا ہوتی ہے:}{\lr{U} \urdu{شکل}}{\lr{V} \urdu{شکل}}{\lr{L} \urdu{شکل}}{\lr{M} \urdu{شکل}}
\umcq{39}{اگر گراف تفاعل کو ظاہر کرے تو ہر عمودی خط کو تفاعل کو ...... پر قطع کرنا چاہیے:}{\lr{4} \urdu{نقاط}}{\lr{3} \urdu{نقاط}}{\lr{2} \urdu{نقاط}}{\lr{1} \urdu{نقطہ}}
\umcq{40}{اگر \lr{$f(x)=x^{3}$} ہو تو \lr{$f(-2)=$}:}{-8}{8}{4}{-610}
\umcq{41}{جملہ \lr{$\dfrac{2x-1}{x^{2}+4}$} ہے:}{\urdu{کثیر رقمی}}{\urdu{الجبری کسر}}{\urdu{عددی کسر}}{\urdu{مساوات}}
\umcq{42}{\lr{$\dfrac{18a^{2}bc^{2}}{12ac}$} کی مختصر ترین شکل ہے:}{\lr{$\dfrac{9abc}{6}$}}{\lr{$\dfrac{3abc}{2}$}}{\lr{$\dfrac{18b}{12}$}}{\lr{$\dfrac{3}{2}$}}
\umcq{43}{\lr{$\dfrac{15xyz}{10}$} کی مختصر ترین شکل ہے:}{\lr{$\dfrac{3xyz}{2}$}}{\lr{$\dfrac{3}{2}$}}{\lr{xyz}}{\lr{$\dfrac{xyz}{2}$}}
\umcq{44}{\lr{$\dfrac{3x+1}{5}+\dfrac{2x+3}{5}$} کی مختصر ترین شکل ہے:}{\lr{$\dfrac{5x+4}{10}$}}{\lr{$\dfrac{5x+4}{5}$}}{\lr{$\dfrac{5x+4}{25}$}}{\lr{5x+4}}
\umcq{45}{\lr{$\left(\dfrac{x^{2}-y^{2}}{x+y}\right)(x-y)$} کی مختصر ترین شکل ہے:}{\lr{$(x+y)^{2}$}}{\lr{$x^{2}+y^{2}$}}{\lr{$x-y$}}{\lr{$(x-y)^{2}$}}
\umcq{46}{\lr{$(x^{3}-y^{3})\div(x-y)$} کی مختصر ترین شکل ہے:}{\lr{$x^{2}-xy+y^{2}$}}{\lr{$x^{2}+xy+y^{2}$}}{\lr{$x^{2}-y^{2}$}}{\lr{$x^{2}+y^{2}$}}
\umcq{47}{مساوات \lr{$\dfrac{5}{x}+\dfrac{1-x}{3}=\dfrac{1}{6x}$} کی قسم ہے:}{\urdu{جذری مساوات}}{\urdu{معکوس مساوات}}{\urdu{کسری مساوات}}{\urdu{قوت نمائی مساوات}}
\umcq{48}{مساوات \lr{$5^{x}+64.5^{-x}-20=0$} کی قسم ہے:}{\urdu{جذری مساوات}}{\urdu{قوت نمائی مساوات}}{\urdu{معکوس مساوات}}{\urdu{کسری مساوات}}
\umcq{49}{\lr{$y^{2}-24y+128=0$} کے روٹس ہیں:}{\lr{8,-16}}{\lr{8,16}}{\lr{-8,16}}{\lr{-8,-16}}
\umcq{50}{\lr{$3x^{2}+10x+3=0$} کے اجزائے ضربی ہیں:}{\lr{$(x+3),(3x+1)$}}{\lr{$(x+3),(3x-1)$}}{\lr{$(x-3),(3x+1)$}}{\lr{$(x-3),(3x-1)$}}
\umcq{51}{\lr{$\cot 60^{o}$} کی قیمت ہے:}{0}{1}{\lr{$\sqrt{3}$}}{\lr{$\dfrac{1}{\sqrt{3}}$}}
\umcq{52}{تمام تکونیاتی نسبتیں مثبت ہوتی ہے:}{\urdu{پہلا ربع}}{\urdu{دوسرا ربع}}{\urdu{تیسرا ربع}}{\urdu{چوتھا ربع}}
\umcq{53}{\lr{$\mathrm{cosec}\,\theta$} مثبت ہوتا ہے:}{\urdu{پہلے اور تیسرے ربع میں}}{\urdu{دوسرے اور تیسرے ربع میں}}{\urdu{پہلے اور دوسرے ربع میں}}{\urdu{پہلے اور چوتھے ربع میں}}
\umcq{54}{سائن کا قانون ہے:}{\lr{$\dfrac{b}{\sin\alpha}=\dfrac{a}{\sin\beta}=\dfrac{c}{\sin\gamma}$}}{\lr{$\dfrac{a}{\sin\alpha}=\dfrac{c}{\sin\beta}=\dfrac{b}{\sin\gamma}$}}{\lr{$\dfrac{a}{\sin\alpha}=\dfrac{b}{\sin\beta}=\dfrac{c}{\sin\gamma}$}}{\lr{$\dfrac{a}{\sin\alpha}=\dfrac{a}{\sin\beta}=\dfrac{a}{\sin\ldots}$}}
\umcq{55}{کوسائن کا قانون ہے:}{\lr{$\cos\alpha=\dfrac{b^{2}+a^{2}-c^{2}}{2ab}$}}{\lr{$\cos\beta=\dfrac{c^{2}+b^{2}-a^{2}}{2cb}$}}{\lr{$\cos\gamma=\dfrac{a^{2}+b^{2}-c^{2}}{2ab}$}}{\lr{$\cos\gamma=\dfrac{ca^{2}+b^{2}-c}{2ab}$}}
\umcq{56}{\lr{$\Delta$ABC} کا رقبہ برابر ہے:}{\lr{$\dfrac{1}{2}ab\sin\alpha\ \dfrac{1}{2}ab\sin\alpha$}}{\lr{$\dfrac{1}{2}bc\sin\beta\ \dfrac{1}{2}bc\sin\beta$}}{\lr{$\dfrac{1}{2}\sin\gamma\ \dfrac{1}{2}\sin\gamma$}}{\lr{$\dfrac{1}{2}ac\sin\beta\ \dfrac{1}{2}ac\sin\beta$}}
\umcq{57}{بیرنگ ...... سے ماپا جاتا ہے:}{\lr{East}}{\lr{West}}{\lr{North}}{\lr{South}}
\umcq{58}{بیرنگ ...... سے لکھا جاتا ہے:}{\lr{figure 1}}{\lr{figures 2}}{\lr{figures 3}}{\lr{figures 4}}
\umcq{59}{\lr{$\tan(180^{o}-\theta)=$}}{\lr{$\tan\theta$}}{\lr{$\cot\theta$}}{\lr{$-\cot\theta$}}{\lr{$-\tan\theta$}}
\umcq{60}{\lr{$\sin(90^{o}+\theta)=$}}{\lr{$\sin\theta$}}{\lr{$-\sin\theta$}}{\lr{$\cos\theta$}}{\lr{$-\cos\theta$}}
\umcq{61}{دائرہ کے محیط پر ایک نقطہ سے مرکز تک فاصلہ کہلاتا ہے:}{\urdu{رداس}}{\urdu{قوس}}{\urdu{وتر}}{\urdu{مماس}}
\umcq{62}{ایک جیسے دائروں کے رداس ہوتے ہیں:}{\urdu{سب نابرابر}}{\urdu{سب برابر}}{\urdu{کسی وتر کا آدھا}}{\urdu{قطر کا دگنا}}
\umcq{63}{دائرہ کی حد کہلاتی ہے:}{\urdu{وتر}}{\urdu{قطعہ}}{\urdu{محیط}}{\urdu{قطر}}
\umcq{64}{محیط کا کوئی حصہ کہلاتا ہے:}{\urdu{وتر}}{\urdu{نصف دائرہ}}{\urdu{رداس}}{\urdu{قوس}}
\umcq{65}{دائرے کے مرکز سے گزرنے والے وتر کو کہتے ہیں:}{\urdu{رداس}}{\urdu{قطر}}{\urdu{قاطع}}{\urdu{محیط}}
\umcq{66}{دی گئی شکل میں \lr{AB} کیا ہے:}{\urdu{قطر}}{\urdu{مماس}}{\urdu{وتر}}{\urdu{قوس}}
\umcq{67}{دی گئی شکل میں قوس کبیرہ ہے:}{\lr{$m\widehat{AB}$}}{\lr{$m\widehat{BC}$}}{\lr{$m\widehat{BDC}$}}{\lr{$m\widehat{AD}$}}
\umcq{68}{تین غیر ہم خط نقاط میں سے ...... دائرہ (دائرے) گزر سکتے ہیں:}{\urdu{ایک}}{\urdu{دو}}{\urdu{تین}}{\urdu{چار}}
\umcq{69}{دائرے کے وتر کا عمودی ناصف دائرہ کے ...... سے گزرتا ہے:}{\urdu{قوس}}{\urdu{رداس}}{\urdu{مرکز}}{\urdu{محیط}}
\umcq{70}{ایک دائرہ کے مرکزی زاویوں کی پیمائشوں کا مجموعہ ہوتا ہے:}{\lr{$90^{o}$}}{\lr{$180^{o}$}}{\lr{$270^{o}$}}{\lr{$360^{o}$}}
\umcq{71}{مماس ایک ایسا خط ہے جو دائرہ کے محیط کو ...... پر مس کرتا ہے:}{\urdu{بہت سے نقاط}}{\urdu{تین نقاط}}{\urdu{دو نقاط}}{\urdu{ایک نقطہ}}
\umcq{72}{دائرے کے بیرونی نقطہ سے دائرہ پر ...... مماس کھینچے جا سکتے ہیں:}{\urdu{ایک}}{\urdu{دو}}{\urdu{تین}}{\urdu{بہت سے}}
\umcq{73}{دائرے کے مرکز سے دائرہ ...... مماس کھینچے جا سکتے ہیں:}{\urdu{صفر}}{\urdu{ایک}}{\urdu{دو}}{\urdu{تین}}
\umcq{74}{دائرہ میں رداس اور مماس ہوتے ہیں:}{\urdu{متوازی}}{\urdu{مساوی}}{\urdu{عمودی}}{\urdu{صفر}}
\umcq{75}{نصف دائرے سے بڑے قطعہ میں بننے والا ...... زاویہ ہوتا ہے:}{\urdu{حادہ}}{\urdu{قائمہ}}{\urdu{منفرجہ}}{\urdu{مستقیم}}
\umcq{76}{دائرے کے مرکز پر قوس سے بننے والا زاویہ کہلاتا ہے:}{\urdu{حادہ زاویہ}}{\urdu{مرکزی زاویہ}}{\urdu{قائمہ زاویہ}}{\urdu{مکمل زاویہ}}
\umcq{77}{نصف دائرہ میں بننے والے زاویہ کی پیمائش کیا ہوتی ہے:}{\lr{$45^{o}$}}{\lr{$60^{o}$}}{\lr{$90^{o}$}}{\lr{$120^{o}$}}
\umcq{78}{دائرے کے ایک ہی قلعہ میں کوئی بھی دو زاویے ہوتے ہیں:}{\urdu{سپلیمنٹری}}{\urdu{کمپلیمنٹری}}{\urdu{مساوی}}{\urdu{صفر}}
\umcq{79}{دائرے کے سیکٹر کا رقبہ = ......:}{\lr{$\dfrac{1}{2}r\theta$}}{\lr{$\dfrac{3}{2}r^{2}\theta$}}{\lr{$\dfrac{1}{2}r^{2}\theta$}}{\lr{$\dfrac{3}{2}r\theta$}}
\umcq{80}{اگر \lr{r=6cm} اور ریڈین \lr{$\theta=2$} ہوں تو قوس کی لمبائی ہے:}{\lr{12 cm}}{\lr{8 cm}}{\lr{6 cm}}{\lr{3 cm}}
\umcq{81}{ملحقہ شکل میں \lr{$\overleftrightarrow{ABC}$} ہے:}{\urdu{ایک وتر}}{\urdu{ایک قوس}}{\urdu{ایک مماس}}{\urdu{ایک قاطع}}
\umcq{82}{قوسیں ...... قسم کی ہوتی ہیں:}{2}{3}{4}{5}
\umcq{83}{دائرے کے وتر کا عمودی ناصف ہمیشہ ...... میں سے گزرتا ہے:}{\urdu{قطر}}{\urdu{تین غیر ہم خط نقاط}}{\urdu{رداس}}{\urdu{مرکز}}
\umcq{84}{دائرے کا صرف ایک وتر ہوتا ہے:}{\urdu{وتر}}{\urdu{مرکز}}{\urdu{قطر}}{\urdu{قاطع}}
\umcq{85}{جہاں دونوں مماس دائرہ کے باہر ملتے ہوں وہ نقطہ بناتا ہے:}{\urdu{نصف دائرہ}}{\urdu{قطر}}{\urdu{رداس}}{\urdu{ایک زاویہ}}
\umcq{86}{ایک نقطہ سے دائرہ پر دو مساوی مماس کھینچے جا سکتے ہیں جب نقطہ ...... ہو:}{\urdu{دائرے کے اوپر}}{\urdu{دائرے کے اندر}}{\urdu{دائرے کے باہر}}{\urdu{دائرے کے مرکز میں}}
\umcq{87}{کسی واحد بیرونی نقطہ سے دائرہ پر کھینچے گئے مماس ہوتے ہیں:}{\urdu{لمبائی میں غیر مساوی}}{\urdu{ایک دوسرے پر عمودی}}{\urdu{لمبائی میں مساوی}}{\urdu{دائرے کے اندر}}
\umcq{88}{دائرہ کا مرکز معلوم کرنے کے لیے کم از کم کتنے وتروں کی ضرورت ہوتی ہے:}{1}{2}{3}{4}
\umcq{89}{ایک قوس کے درمیانی نقطہ پر مماس کھینچنے کا پہلا مرحلہ ہوتا ہے:}{\urdu{قوس کے اختتامی نقاط کو ملانا}}{\urdu{رداس کھینچنا}}{\urdu{وتر پر عمودی لائن کھینچنا}}{\urdu{ایک زاویہ کھینچنا}}
\umcq{90}{رداس اور مماس کے درمیان نقطہ تماس پر زاویہ ہوتا ہے:}{\lr{$30^{o}$}}{\lr{$60^{o}$}}{\lr{$90^{o}$}}{\lr{$120^{o}$}}
\umcq{91}{مجموعی تعدادات معلوم کرنے کے لیے استعمال ہوتا ہے:}{\urdu{جمع}}{\urdu{جذر المربع}}{\urdu{ضرب}}{\urdu{تقسیم}}
\umcq{92}{دوسرا چہارمی حصہ ظاہر کرتا ہے:}{\urdu{اوسط}}{\urdu{عادہ}}{\urdu{وسطانیہ}}{\urdu{تغیریت}}
\umcq{93}{پہلا چہارمی حصہ مواد کو ....... برابر حصوں میں تقسیم کرتا ہے:}{\urdu{دو}}{\urdu{تین}}{\urdu{چار}}{\urdu{دس}}
\umcq{94}{سب سے بڑی اور سب سے چھوٹی قدر کے درمیان فرق کہلاتا ہے:}{\urdu{اوسط}}{\urdu{تغیریت}}{\urdu{سعت}}{\urdu{معیاری انحراف}}
\umcq{95}{انتشاری خاکہ ...... متغیرات کے درمیان تعلق کو ظاہر کرتا ہے:}{\urdu{پانچ}}{\urdu{دو}}{\urdu{تین}}{\urdu{چار}}
\umcq{96}{انتشار کا پیمانہ ہے:}{\urdu{اوسط}}{\urdu{وسطانیہ}}{\urdu{عادہ}}{\urdu{تغیریت}}
\umcq{97}{تغیریت کا مثبت جذر المربع کہلاتا ہے:}{\urdu{اوسط}}{\urdu{وسطانیہ}}{\urdu{معیاری انحراف}}{\urdu{سعت}}
\umcq{98}{...... انتشار کا پیمانہ نہیں ہے:}{\urdu{سعت}}{\urdu{حسابی اوسط}}{\urdu{تغیریت}}{\urdu{معیاری انحراف}}
\umcq{99}{مواد \lr{8,8,8,8,8,8} کی تغیریت ہے:}{0}{16}{8}{48}
\umcq{100}{پہلے \lr{20} قدرتی اعداد کی سعت ہے:}{20}{10}{19}{30}
\umcq{101}{ناممکن واقعہ کا احتمال ہے:}{0}{1}{2}{-1}
\umcq{102}{ایک سکے کو ایک بار اچھالنے پر ہیڈ حاصل کرنے کا احتمال کیا ہے:}{0}{0.25}{0.5}{1}
\umcq{103}{کسی تجربے میں تمام ممکنہ نتائج کے احتمالات کا مجموعہ ہوتا ہے:}{0.5}{0.25}{1}{0.4}
\umcq{104}{اگر \lr{P(A)=0.6} ہو تو واقعہ \lr{A} وقوع پذیر نہ ہونے کا احتمال ہے:}{0.4}{0.6}{1}{1.6}
\umcq{105}{دو واقعات باہمی متناہی کہلاتے ہیں اگر:}{\urdu{وہ ایک ہی وقت میں وقوع پذیر ہو سکتے ہوں}}{\urdu{ایک دوسرے کو متاثر کرتے ہوں}}{\urdu{وہ ایک ساتھ وقوع پذیر نہ ہو سکتے ہوں}}{\urdu{وہ ہمیشہ برابر ہوں}}
\umcq{106}{اگر ایک سکہ اچھالا جائے اور ایک ڈائس کو پھینکا جائے تو سیمپل پوائنٹس کی تعداد ہے:}{3}{6}{2}{12}
\umcq{107}{یقینی واقعہ کا احتمال کیا ہوتا ہے:}{0}{1}{\lr{$\dfrac{2}{3}$}}{\lr{$\dfrac{4}{5}$}}
\umcq{108}{ایک ڈائس پر \lr{6} سے بڑا عدد حاصل کرنے کا احتمال ہے:}{0.33}{1}{0}{0.5}
\umcq{109}{تاش کی گڈی کے پتوں میں سے ایک پتا نکالنے کے ...... ممکنہ نتائج ہوتے ہیں:}{13}{1}{52}{26}
\umcq{110}{اگر \lr{P(E)=0.07} ہو تو \lr{E} نہ ہونے کا احتمال ہے:}{0.95}{0.89}{0.93}{0.90}

\end{document}
